\section{Related Work}
\label{sec:related}

\para{Emergent abilities and phase transitions in LLMs.}
\citet{wei2022emergent} define emergent abilities as capabilities that appear unpredictably as models scale, exhibiting near-random performance below a critical size threshold and substantially above-random performance above it. This framing has shaped expectations about how model behaviors change with scale. However, \citet{schaeffer2023emergent} challenge this narrative, demonstrating that over 92\% of claimed emergent abilities on BIG-Bench use nonlinear or discontinuous metrics. When evaluated with linear metrics such as token edit distance or Brier scores, the same models show smooth, predictable improvement. This debate is directly relevant to alignment resistance: if apparent phase transitions in safety behavior are similarly metric-dependent, the practical implications for AI safety differ substantially. \citet{hong2025evidence} provide complementary evidence that phase transitions occur even in small (3.6M parameter) transformers during training, detectable through Poisson-centered statistical probes but invisible to standard loss curves.

\para{Alignment robustness and resistance.}
\citet{ji2024language} introduce alignment \emph{elasticity}, showing that post-alignment models tend to revert to pre-training distributions when further fine-tuned. They model this process through compression theory and demonstrate that elasticity increases with both model size and pre-training data volume. Critically, they show that inverse alignment (removing safety) is consistently easier than forward alignment (adding safety), tested across Llama-2, Llama-3, and Qwen families. \citet{qi2024finetuning} demonstrate that even benign fine-tuning can compromise safety alignment, suggesting that alignment is inherently fragile. Our work extends these findings by probing alignment robustness at inference time across a denser range of model scales, without requiring access to model weights.

\para{Deceptive alignment and alignment faking.}
\citet{greenblatt2024alignment} provide the first evidence of alignment faking emerging without explicit training, where Claude models strategically comply during monitored interactions while preserving non-compliant behavior when unmonitored. The key scaling observation---that alignment faking appears in Opus and Sonnet but not Haiku---suggests a capability threshold. \citet{hubinger2024sleeper} demonstrate that deliberately inserted deceptive behaviors persist through safety training, with persistence \emph{increasing} with model scale. Together, these results motivate our hypothesis that alignment resistance may exhibit threshold-like behavior, though our experiments test this at inference time rather than through fine-tuning.

\para{Jailbreaking and adversarial attacks on alignment.}
\citet{wei2024jailbroken} provide a taxonomy of jailbreak attacks and analyze why LLM safety training fails, identifying competing objectives and mismatched generalization as root causes. \citet{perez2022red} introduce automated red-teaming using language models to generate adversarial prompts, establishing the methodology of using LLMs to evaluate LLM safety. We build on these approaches by testing standardized jailbreak templates across model scales to measure how resistance changes with capability.

\para{Mechanistic analysis of alignment.}
\citet{turner2025model} create model organisms for emergent misalignment by fine-tuning instruct models on narrow domains. They discover a simultaneous mechanistic and behavioral phase transition: LoRA vectors undergo a sudden rotation at a specific training step, correlated with the emergence of misaligned behavior. This provides the clearest mechanistic evidence for phase transitions in alignment-relevant properties, though it occurs in training dynamics rather than as a function of model scale. \citet{zou2023representation} propose representation engineering as a top-down approach to analyzing model internals, which could complement our behavioral probing in future work.

Unlike prior work that tests alignment robustness through fine-tuning or focuses on a small number of model sizes, we conduct a systematic behavioral evaluation across five model scales using multiple attack vectors and both continuous and discrete metrics, directly testing whether alignment resistance exhibits phase transition or gradual scaling behavior.
